%% ma: L10-B02-0018
%% lop: 10
%% chuong: 1
%% bai: 2
%% dang: 10.2.10
%% loai: DS
%% muc: [TH, TH, TH, TH]
%% dap_an: "ĐSĐS"
%% nguon: "Ảnh giáo viên gửi 10/10/2026 (ThS. Nguyễn Hoàng Việt, Chương 2 Tập hợp), A-Đề 1, Phần 2 Câu 16"
%% kien_thuc: "Bài 2 (tập hợp và các phép toán trên tập hợp)"
%% trang_thai: cho-duyet
%% ly_do: "Dạng toán mới đề xuất 10.2.10: hai tập hợp bằng nhau"
%% ghi_chu: "Đề gốc ghi như một mệnh đề điều kiện (A=B=C, khi đó các ý a-d); hiểu mỗi ý là kiểm tra cặp (x,y) cho có thỏa A=B=C không"
\begin{ex}
Cho ba tập hợp $A=\{2;5\}$, $B=\{5;x\}$, $C=\{x;y;5\}$, biết $A=B=C$. Khi đó
\choiceTF
{\True $x=y=2$ thì $A=B=C$.}
{$x=y=3$ thì $A=B=C$.}
{\True $x=2,\ y=5$ thì $A=B=C$.}
{$x=1,\ y=3$ thì $A=B=C$.}
\loigiai{a) $B=\{5;2\}=A$, $C=\{2;2;5\}=\{2;5\}=A$. Đúng. b) $B=\{5;3\}\ne A$. Sai. c) $B=\{5;2\}=A$, $C=\{2;5;5\}=\{2;5\}=A$. Đúng. d) $B=\{5;1\}\ne A$. Sai.}
\end{ex}
